\documentclass[12pt]{article}
\usepackage{../../ficha}

\begin{document}

\ficha{Derivada, Reta Tangente e Regras de Derivação}

\instrucoes{Resolve os exercícios, apresentando os cálculos. A derivada de $f$ no ponto $a$ é $f'(a) = \displaystyle\lim_{h \to 0} \dfrac{f(a + h) - f(a)}{h}$.}

\secao{Taxa Média de Variação}{Calcula a taxa média de variação de $f$ no intervalo indicado.}
\begin{exercicios}
  \item $f(x) = 3x + 2$, \; em $[1, 4]$
  \item $f(x) = x^2$, \; em $[1, 3]$
  \item $f(x) = x^2 - 4x$, \; em $[0, 2]$
  \item $f(x) = \dfrac{1}{x}$, \; em $[1, 2]$
  \item $f(x) = \sqrt{x}$, \; em $[4, 9]$
  \item A posição de um carro, em km, é $s(t) = 60t + 5t^2$ ($t$ em horas). Calcula a velocidade média entre $t = 1$ e $t = 3$.
  \item Interpreta geometricamente a taxa média de variação.
\end{exercicios}

\secao{Derivada por Definição}{Calcula $f'(a)$ usando a definição.}
\begin{exercicios}
  \item $f(x) = 5x - 1$, \; $a = 2$
  \item $f(x) = x^2$, \; $a = 3$
  \item $f(x) = 2x^2 + x$, \; $a = 1$
  \item $f(x) = x^3$, \; $a = 1$
  \item $f(x) = \dfrac{1}{x}$, \; $a = 2$
  \item $f(x) = \sqrt{x}$, \; $a = 4$
  \item Mostra, pela definição, que a derivada de $f(x) = x^2$ é $f'(x) = 2x$.
  \item Mostra que $f(x) = |x|$ não tem derivada em $x = 0$.
\end{exercicios}

\secao{Regras de Derivação}{Determina a função derivada.}
\begin{exercicios}
  \item $f(x) = 4x^3 - 2x^2 + 7x - 5$
  \item $f(x) = \dfrac{x^4}{2} - \dfrac{3}{x}$
  \item $f(x) = 6\sqrt{x} + \dfrac{1}{x^2}$
  \item $f(x) = (2x + 1)(x^2 - 3)$
  \item $f(x) = x^2\sqrt{x}$
  \item $f(x) = \dfrac{x + 1}{x - 1}$
  \item $f(x) = \dfrac{3x}{x^2 + 1}$
  \item $f(x) = \dfrac{x^2}{2x - 1}$
  \item $f(x) = (3x - 2)^5$
  \item $f(x) = \sqrt{x^2 + 4}$
  \item $f(x) = \dfrac{1}{(x^2 + 1)^3}$
  \item $f(x) = x\,(1 - x)^4$
\end{exercicios}

\secao{Reta Tangente}{Determina a equação reduzida da reta tangente ao gráfico de $f$ no ponto de abcissa indicada.}
\begin{exercicios}
  \item $f(x) = x^2 - 3x$, \; $x = 2$
  \item $f(x) = x^3 - 2x + 1$, \; $x = 1$
  \item $f(x) = \dfrac{4}{x}$, \; $x = 2$
  \item $f(x) = \sqrt{x}$, \; $x = 9$
  \item $f(x) = \dfrac{x}{x + 1}$, \; $x = 0$
  \item Determina o ponto do gráfico de $f(x) = x^2 - 6x + 10$ onde a reta tangente é horizontal.
  \item Determina os pontos do gráfico de $f(x) = x^3 - 3x$ onde a reta tangente é paralela à reta $y = 9x$.
  \item Determina a reta tangente a $f(x) = x^2$ que é paralela à reta $y = 4x - 1$.
\end{exercicios}

\secao{Velocidade e Aceleração}{Resolve.}
\begin{exercicios}
  \item A posição de uma partícula é $s(t) = t^3 - 6t^2 + 9t$ (em metros, $t$ em segundos). Determina a velocidade e a aceleração no instante $t$.
  \item Na mesma partícula, em que instantes a velocidade é nula?
  \item Uma bola é lançada verticalmente e a altura é $h(t) = 30t - 5t^2$. Qual é a velocidade ao fim de 2\,s? Quando atinge a altura máxima?
  \item Determina a segunda derivada de $f(x) = x^4 - 2x^3 + x$.
\end{exercicios}

\secao{Desafios}{Resolve os problemas.}
\begin{exercicios}
  \item Determina $a$ e $b$ de modo que a reta tangente ao gráfico de $f(x) = ax^2 + bx$ no ponto $(1, 5)$ tenha declive 8.
  \item Determina $k$ de modo que a reta $y = 2x + k$ seja tangente ao gráfico de $f(x) = x^2 + 1$.
  \item Mostra que as retas tangentes a $f(x) = \dfrac{1}{x}$ formam com os eixos coordenados triângulos de área constante. Qual é essa área?
  \item Determina $a$ e $b$ de modo que $f(x) = \begin{cases} x^2 & \text{se } x \le 1 \\ ax + b & \text{se } x > 1 \end{cases}$ seja diferenciável em $x = 1$.
\end{exercicios}

\end{document}
